% Options for packages loaded elsewhere
\PassOptionsToPackage{unicode}{hyperref}
\PassOptionsToPackage{hyphens}{url}
\documentclass[
  ignorenonframetext,
]{beamer}
\newif\ifbibliography
\usepackage{pgfpages}
\setbeamertemplate{caption}[numbered]
\setbeamertemplate{caption label separator}{: }
\setbeamercolor{caption name}{fg=normal text.fg}
\beamertemplatenavigationsymbolsempty
% remove section numbering
\setbeamertemplate{part page}{
  \centering
  \begin{beamercolorbox}[sep=16pt,center]{part title}
    \usebeamerfont{part title}\insertpart\par
  \end{beamercolorbox}
}
\setbeamertemplate{section page}{
  \centering
  \begin{beamercolorbox}[sep=12pt,center]{section title}
    \usebeamerfont{section title}\insertsection\par
  \end{beamercolorbox}
}
\setbeamertemplate{subsection page}{
  \centering
  \begin{beamercolorbox}[sep=8pt,center]{subsection title}
    \usebeamerfont{subsection title}\insertsubsection\par
  \end{beamercolorbox}
}
% Prevent slide breaks in the middle of a paragraph
\widowpenalties 1 10000
\raggedbottom
\AtBeginPart{
  \frame{\partpage}
}
\AtBeginSection{
  \ifbibliography
  \else
    \frame{\sectionpage}
  \fi
}
\AtBeginSubsection{
  \frame{\subsectionpage}
}
\usepackage{iftex}
\ifPDFTeX
  \usepackage[T1]{fontenc}
  \usepackage[utf8]{inputenc}
  \usepackage{textcomp} % provide euro and other symbols
\else % if luatex or xetex
  \usepackage{unicode-math} % this also loads fontspec
  \defaultfontfeatures{Scale=MatchLowercase}
  \defaultfontfeatures[\rmfamily]{Ligatures=TeX,Scale=1}
\fi
\usepackage{lmodern}
\ifPDFTeX\else
  % xetex/luatex font selection
\fi
% Use upquote if available, for straight quotes in verbatim environments
\IfFileExists{upquote.sty}{\usepackage{upquote}}{}
\IfFileExists{microtype.sty}{% use microtype if available
  \usepackage[]{microtype}
  \UseMicrotypeSet[protrusion]{basicmath} % disable protrusion for tt fonts
}{}
\makeatletter
\@ifundefined{KOMAClassName}{% if non-KOMA class
  \IfFileExists{parskip.sty}{%
    \usepackage{parskip}
  }{% else
    \setlength{\parindent}{0pt}
    \setlength{\parskip}{6pt plus 2pt minus 1pt}}
}{% if KOMA class
  \KOMAoptions{parskip=half}}
\makeatother
\setlength{\emergencystretch}{3em} % prevent overfull lines
\providecommand{\tightlist}{%
  \setlength{\itemsep}{0pt}\setlength{\parskip}{0pt}}
\usepackage{bookmark}
\IfFileExists{xurl.sty}{\usepackage{xurl}}{} % add URL line breaks if available
\urlstyle{same}
\hypersetup{
  hidelinks,
  pdfcreator={LaTeX via pandoc}}

\author{\texorpdfstring{}{}}
\date{}

\begin{document}

\begin{frame}[fragile]{Security and Provenance}
\protect\phantomsection\label{security-and-provenance}
Research integrity requires more than reproducibility; it requires
verifiable authorship. In an era of generative AI, automated scraping,
and synthetic media, the ability to prove that a document was produced
by a specific pipeline at a specific time is a critical defense against
fabrication and misattribution. The W3C PROV data model
{[}@moreau2013provdm{]} establishes a formal vocabulary for expressing
provenance records---entities, activities, and agents connected by
derivation, generation, and attribution relations. The Docxology
Template implements a domain-specific provenance layer that embeds these
relations directly into the PDF artifact itself, using four
complementary steganographic and cryptographic mechanisms.

\begin{block}{Threat Model}
\protect\phantomsection\label{threat-model}
The steganography subsystem defends against three classes of threats:

\begin{enumerate}
\tightlist
\item
  \textbf{Unauthorized redistribution}: A manuscript is scraped and
  republished without attribution. The steganographic watermark survives
  the redistribution and can be used to prove original authorship.
\item
  \textbf{Content tampering}: Figures or text are modified after
  publication. The SHA-256 hash embedded in the PDF metadata detects any
  modification to the file contents.
\item
  \textbf{Provenance forgery}: An attacker claims to have produced a
  document they did not author. The build timestamp, commit hash, and
  pipeline metadata embedded in the watermark provide a verifiable chain
  of custody.
\end{enumerate}
\end{block}

\begin{block}{Steganographic Layers}
\protect\phantomsection\label{steganographic-layers}
The system applies four complementary layers of provenance information:

\begin{block}{Layer 1: PDF Metadata Injection}
\protect\phantomsection\label{layer-1-pdf-metadata-injection}
The \texttt{inject\_pdf\_metadata} function writes structured metadata
into both the PDF Info dictionary and an XMP (Extensible Metadata
Platform) packet:

\begin{itemize}
\tightlist
\item
  \texttt{/Creator}: Pipeline identifier
\item
  \texttt{/Producer}: Module path
  (\texttt{infrastructure.steganography})
\item
  \texttt{/CreationDate}: UTC timestamp in ISO 8601 format
\item
  \texttt{/Author}: From \texttt{config.yaml}
\item
  \texttt{/Title}: From \texttt{config.yaml}
\item
  Custom fields: DOI, ORCID, repository URL
\end{itemize}
\end{block}

\begin{block}{Layer 2: Cryptographic Hashing}
\protect\phantomsection\label{layer-2-cryptographic-hashing}
Before watermarking, a SHA-256 hash of the rendered PDF is computed and
stored in:

\begin{itemize}
\tightlist
\item
  The output manifest (\texttt{output/manifest.json})
\item
  The PDF metadata (\texttt{/Subject} field)
\item
  An external hash file
  (\texttt{output/\textless{}name\textgreater{}.sha256})
\end{itemize}

This enables post-hoc verification: anyone with the hash can verify that
the PDF has not been modified since rendering.
\end{block}

\begin{block}{Layer 3: Alpha-Channel Text Overlay}
\protect\phantomsection\label{layer-3-alpha-channel-text-overlay}
A semi-transparent text overlay is applied to each page of the PDF,
encoding:

\begin{itemize}
\tightlist
\item
  Build timestamp
\item
  Git commit hash (short SHA)
\item
  Project name
\item
  Pipeline version
\end{itemize}

The overlay is rendered at low opacity (typically 3--5\% alpha) to be
invisible during normal viewing but detectable through image analysis.
It survives printing (as a faint watermark) and standard PDF operations.
\end{block}

\begin{block}{Layer 4: QR Code Injection}
\protect\phantomsection\label{layer-4-qr-code-injection}
An optional QR code is generated containing a URL pointing to the
repository (e.g., \texttt{github.com/docxology/template}). The QR code
is placed in a configurable position (default: bottom-right corner of
the last page) at a specified size.
\end{block}
\end{block}

\begin{block}{The \texttt{secure\_run.sh} Orchestrator}
\protect\phantomsection\label{the-secure_run.sh-orchestrator}
The steganographic pipeline is orchestrated by \texttt{secure\_run.sh},
a Bash script that wraps the standard \texttt{run.sh} pipeline with
post-processing steganography:

\begin{enumerate}
\tightlist
\item
  Execute the standard 8-stage pipeline for the target project.
\item
  Locate the rendered PDF in \texttt{output/}.
\item
  Apply metadata injection, hashing, text overlay, and QR code
  injection.
\item
  Save the secured PDF alongside the original.
\item
  Generate a steganography report in JSON format.
\end{enumerate}

The orchestrator processes either a single specified project or all
discovered projects sequentially.
\end{block}

\begin{block}{Tamper Detection}
\protect\phantomsection\label{tamper-detection}
Verification is performed by comparing the stored SHA-256 hash against a
freshly computed hash of the distributed PDF. Any modification---even a
single bit flip---produces a hash mismatch. The alpha-channel overlay
provides a secondary, visual verification channel that does not require
access to the original hash.
\end{block}

\begin{block}{Limitations}
\protect\phantomsection\label{limitations}
\begin{itemize}
\tightlist
\item
  Alpha-channel overlays are stripped by some PDF optimization tools
  (e.g., \texttt{qpdf\ -\/-optimize}).
\item
  QR codes are visible and may be removed by a determined attacker.
\item
  The system does not provide non-repudiation in the cryptographic
  sense---it does not use digital signatures with private keys. Future
  versions may integrate GPG or X.509 signing.
\item
  The provenance model is pipeline-centric rather than PROV-compliant;
  future work will generate machine-readable PROV-O
  {[}@moreau2013provdm{]} records alongside the embedded watermarks.
\end{itemize}
\end{block}

\begin{block}{Relationship to Software Supply Chain Integrity}
\protect\phantomsection\label{relationship-to-software-supply-chain-integrity}
The steganographic provenance layer operates at the \emph{document
level}---it certifies the integrity of a specific PDF artifact. A
complementary concern is \emph{build-level} provenance: certifying that
the pipeline itself was executed with verified source code and
dependencies. Frameworks such as in-toto {[}@torresarias2019intoto{]}
and SLSA (Supply-chain Levels for Software Artifacts) address this
concern by defining attestation chains from source commit through build
steps to final artifact. Future versions of the Docxology Template may
generate SLSA-compatible provenance attestations alongside the
steganographic watermarks, creating a two-layer provenance model:
in-toto attests that the build pipeline was executed with the claimed
source code, while the steganographic layer attests that the PDF was
produced by that pipeline at a specific time.
\end{block}

\begin{block}{Relationship to FAIR and Formal Provenance Standards}
\protect\phantomsection\label{relationship-to-fair-and-formal-provenance-standards}
The steganographic layer supports the FAIR for Research Software
(FAIR4RS) principles {[}@barker2022fair4rs{]} at the artifact level.
PDFs carry embedded metadata (Findability) in standardized XMP format
(Interoperability). The SHA-256 hash manifest enables persistent
integrity verification (a prerequisite for Reusability). The
Documentation Duality standard ensures that the software producing the
artifact is inspectable and well-documented (satisfying FAIRsoft
{[}@garijo2024fairsoft{]} metadata and documentation indicators). Full
PROV-compliant provenance traces---capturing the derivation chain from
source data through analysis scripts to rendered PDF---are a natural
extension and a primary target for future development.
\end{block}
\end{frame}

\end{document}
